\documentclass[11pt]{article}
\usepackage{amsmath}
\title{Changing a lift session does not retire an old payload effect}
\date{26 November 2026}
\begin{document}
\maketitle

\paragraph{Source cut.}
This fixture executes the pinned Open-RMF supervisor request and
state methods, together with the corresponding simulation plugin's
request body, state construction, geometric payload filter and
per-payload velocity write. Minimal stand-ins supply ROS publications,
poses, component storage, and a selected actual nonzero cabin-velocity
input. The fixture explicitly chooses when the floor-and-door command
finishes. It does not run the Gazebo scheduler, decide when a real
cabin can move, or run the mission planner. The fleet end-session
phase is inspected, not executed: its completion predicate compares
the reported session with its requester. In normal mission context
other phases may ensure egress before this phase is called.

\paragraph{Native protocol reconciliation.}
With a prior command active, an end-session request is different
from the simulator's current command and the driver declines it.
The supervisor has already removed its corresponding active pointer
when publishing the end, so it accepts the other fleet's next
request; the busy simulator declines that request too. Once
the floor-and-door command has completed, a report of the previous
session prompts the supervisor, which now has the successor active,
to reissue the successor request. The simulator accepts that
request at the idle command cut. This is ordinary modeled event
order, not a transport, security or protocol attack. A second trace
where the first end command was accepted also yields the successor
session when it is next admitted. With and without an accepted END,
the selected callbacks reconcile state toward the same successor.
The simulated busy-command gate is a real distinction and is not
removed from the comparator.

\paragraph{Physical effect scope.}
The plugin takes payloads from model positions in the translated
initial cabin bounding box, not from the fleet session identifier.
For models with a linear-velocity command component, the positive
cabin-velocity branch writes its vertical command to every selected
payload. In the matched executions the same lift reports and command
cuts coexist with the old model either inside or outside. Only the
inside model receives the successor move's velocity write. This is
also possible after a genuinely accepted END; no session
acknowledgement is a body-exit test. Even a clear payload lookup
at successor acceptance fails to settle the future effect if a model
enters between that check and the later payload lookup. The replay
adds exactly that difference in positions while keeping the admitted
command and reports unchanged. The unrelated model outside the
cabin receives no command in any case.

\paragraph{Result and limits.}
Let $C$ be a modeled command and $A(C)$ the set of models whose
receipt of its vertical effect is authorized. At a nonzero velocity
update, let $M$ be the models selected by the actual payload filter
and with the necessary velocity component. The per-effect inclusion
law is $M\subseteq A(C)$, not merely acceptance of $C$ or the end of
the old session. The native effect loop makes this condition exact
for that modeled write: each selected model is written, and the
loop does not write to others. This does not prove future motion
or that a velocity command is physically obeyed.
The matched witness proves the lift session and state projection
cannot decide this condition, even with the busy gate and nominal
accepted-END trace. A geometric witness at issue also cannot decide
it without containment through the actual use cut. An effect-side
authorization source, physical exclusion covering the interval,
or an explicated recheck can close the modeled gap; no such
mechanism is attributed to the stock simulator here.

The box uses model reference positions, not a complete cabin
occupancy or clearance sensor and not a tail/body certificate.
There is no claim of field behavior, lift safety, authorized
human transport, operating service bound, full R5 or architecture
acceptance. This is a bounded source-backed transfer of the
issued-to-effect distinction from roadway and door fixtures to
a platform that can carry a previous participant.
\end{document}
